\chapter{20 तक गिनती}\label{ch:g1:counting}

कितने कंचे हैं? गिनती गणित के सबसे पहले प्रश्न का उत्तर देती है। इस
अध्याय में हम छोटे-छोटे ढेर गिनेंगे, $0$ से $20$ तक की संख्याएँ
लिखेंगे, और तुलना करेंगे: किसके पास ज़्यादा हैं?

\section{वस्तुएँ गिनना}

\begin{method}[गिनने का तरीका]\label{met:g1:counting:how}
\begin{enumerate}
  \item हर वस्तु को एक बार छुओ और क्रम से बोलते जाओ: एक,
        दो, तीन, \dots;
  \item कोई वस्तु न छूटे, कोई दो बार न छुए (गिनी हुई वस्तुएँ
        अलग रखते जाओ, या उन पर निशान लगाते जाओ);
  \item जो \emph{आख़िरी} संख्या बोली गई, उतनी ही वस्तुएँ हैं।
\end{enumerate}
छूने का क्रम कोई मायने नहीं रखता: बाएँ से दाएँ गिनो या
दाएँ से बाएँ, $7$ कंचे $7$ ही रहते हैं।
\end{method}

\begin{omfigure}
\begin{tikzpicture}[scale=0.75]
  \foreach \i/\n in {0/1,1/2,2/3,3/4,4/5} {
    \fill[omDef] (\i*1.1,0) circle (0.33);
    \node[above, font=\small] at (\i*1.1,0.45) {\n};
  }
  \node[right, font=\small] at (5.2,0) {पाँच कंचे: कुल $5$};
\end{tikzpicture}

{\small गिनती: हर कंचे के लिए एक संख्या, और आख़िरी संख्या ही
उत्तर है।}
\end{omfigure}

\begin{definition}[शून्य]\label{def:g1:counting:zero}
जब गिनने के लिए कुछ भी न हो --- एक ख़ाली डिब्बा --- तब संख्या
\emph{शून्य}\index{शून्य} होती है, जिसे $0$ लिखते हैं।
\end{definition}

\section{20 तक की संख्याएँ}

\begin{example}[पढ़ना और लिखना]\label{ex:g1:counting:writing}
$0$ से $20$ तक की संख्याएँ, क्रम से:
\[
0,\ 1,\ 2,\ 3,\ 4,\ 5,\ 6,\ 7,\ 8,\ 9,\ 10,
\]
\[
11,\ 12,\ 13,\ 14,\ 15,\ 16,\ 17,\ 18,\ 19,\ 20 .
\]
$9$ के बाद संख्याएँ \emph{दो} अंकों की हो जाती हैं: $13$ यानी एक
दहाई और $3$ इकाई; $20$ यानी दो दहाई। (दहाइयों की बारी
\cref{ch:g1:tens} में आती है।)
\end{example}

\begin{omfigure}
\begin{tikzpicture}[scale=0.62]
  \draw[-Stealth] (-0.4,0) -- (20.8,0);
  \foreach \i in {0,...,20} {
    \draw (\i,-0.12) -- (\i,0.12);
    \node[below, font=\footnotesize] at (\i,-0.15) {\i};
  }
  \fill[omThm] (7,0) circle (2.6pt);
  \fill[omMeth] (13,0) circle (2.6pt);
\end{tikzpicture}

{\small $0$ से $20$ तक की संख्या रेखा: हर संख्या का अपना स्थान है।
दाईं ओर बढ़ने पर संख्याएँ बड़ी होती जाती हैं। बिंदु $7$ और $13$ पर हैं।}
\end{omfigure}

\begin{example}[ठीक पहले, ठीक बाद]\label{ex:g1:counting:neighbors}
रेखा पर हर संख्या के दो पड़ोसी होते हैं: $14$ से ठीक पहले $13$ आता
है, ठीक बाद $15$। उलटी गिनती --- $10$, $9$, $8$,
\dots\ --- बाईं ओर चलना है।
\end{example}

\section{तुलना करना}

\begin{method}[किसके पास ज़्यादा हैं?]\label{met:g1:counting:compare}
दो ढेरों की तुलना करने के लिए जोड़े बनाओ: हर ढेर से एक-एक,
आमने-सामने।
\begin{enumerate}
  \item अगर एक ढेर की कुछ वस्तुएँ बिना जोड़ीदार के बच जाएँ, तो
        उस ढेर में \emph{ज़्यादा} हैं;
  \item अगर दोनों ओर कुछ न बचे, तो दोनों में \emph{बराबर}
        हैं।
\end{enumerate}
संख्याओं के साथ जोड़े बनाने की ज़रूरत ही नहीं: संख्या रेखा पर जो
दाईं ओर है वही बड़ी है --- जोड़े और रेखा हमेशा एक ही बात कहते हैं।
हम लिखते हैं $12 > 9$ ($12$ ज़्यादा है) या $9 < 12$ ($9$ कम है)।
\end{method}

\begin{omfigure}
\begin{tikzpicture}[scale=0.62]
  \foreach \i in {0,...,6} \fill[omDef] (\i*1.1,1.3) circle (0.3);
  \foreach \i in {0,...,4} \fill[omThm] (\i*1.1,0) circle (0.3);
  \foreach \i in {0,...,4}
    \draw[omExo, thick] (\i*1.1,0.3) -- (\i*1.1,1.0);
  \draw[omExo, thick] (5.5,1.3) circle (0.52);
  \draw[omExo, thick] (6.6,1.3) circle (0.52);
  \node[right, font=\small] at (7.5,0.65)
    {$7 > 5$: दो नीले कंचों का जोड़ीदार नहीं};
\end{tikzpicture}

{\small जोड़े बनाकर तुलना: $7$ नीले कंचों के सामने $5$ लाल
--- दो नीले (घेरे हुए) बच जाते हैं, यानी $7$, $5$ से ज़्यादा है।}
\end{omfigure}

\begin{example}[पहला, दूसरा, तीसरा]\label{ex:g1:counting:ordinal}
दौड़ में धावक क्रम से पहुँचते हैं: \emph{पहला},
\emph{दूसरा}, \emph{तीसरा}, \emph{चौथा}\,\dots\ ये
शब्द \emph{स्थान} बताते हैं, गिनती नहीं: तीसरा धावक एक ही धावक
है, जो स्थान $3$ पर आया।
\end{example}

\section{अभ्यास}

\begin{exercise}[$\star$]\label{exo:g1:counting:1}
अपने नाम के अक्षर गिनो। दो हाथों की उँगलियाँ गिनो। दो बिल्लियों
के पैर गिनो।
\end{exercise}

\begin{exercise}[$\star$]\label{exo:g1:counting:2}
अंकों में लिखो: सात; \ बारह; \ बीस; \ \omterm{def:g1:counting:zero}{शून्य}। शब्दों में लिखो:
$4$; \ $15$; \ $18$।
\end{exercise}

\begin{exercise}[$\star$]\label{exo:g1:counting:3}
गिनती आगे बढ़ाओ: $8,\ 9,\ 10,\ ?,\ ?,\ ?$। उलटी गिनती करो:
$15,\ 14,\ 13,\ ?,\ ?,\ ?$।
\end{exercise}

\begin{exercise}[$\star$]\label{exo:g1:counting:4}
$10$ से ठीक पहले कौन-सी संख्या आती है? $16$ के ठीक बाद? $12$
और $14$ के बीच?
\end{exercise}

\begin{exercise}[$\star$]\label{exo:g1:counting:5}
उतार कर $<$ या $>$ से पूरा करो:
\[
6 \;?\; 9, \qquad
14 \;?\; 8, \qquad
11 \;?\; 13, \qquad
20 \;?\; 19 .
\]
\end{exercise}

\begin{exercise}[$\star$]\label{exo:g1:counting:6}
$6$ वृत्त और $8$ क्रॉस बनाओ। उनके जोड़े बनाओ: कौन ज़्यादा हैं, और
कितने से?
\end{exercise}

\begin{exercise}[$\star$]\label{exo:g1:counting:7}
सबसे छोटी से सबसे बड़ी के क्रम में लिखो: $12$; \ $7$; \ $17$; \ $2$; \ $10$।
\end{exercise}

\begin{exercise}[$\star$]\label{exo:g1:counting:8}
पाँच बच्चे पंक्ति में खड़े हैं: अना, बॉब, क्लियो, डैन, एम्मा। दूसरे
स्थान पर कौन है? चौथे पर कौन? डैन किस स्थान पर है?
\end{exercise}

\begin{exercise}[$\star\star$]\label{exo:g1:counting:9}
मैं $10$ और $20$ के बीच की एक संख्या हूँ। रेखा पर गिनने से मैं
$13$ के तीन कदम बाद पड़ती हूँ। मैं कौन हूँ? और $13$ से तीन कदम
\emph{पहले} कौन है?
\end{exercise}
